These are the bounds the question asks for in its strictest form: a deterministic bound on the error of the trained network itself, computed from the training set alone, with no validation data, no unlabeled data and no knowledge of the input distribution.

\paragraph{Worst-case weight sets (\cref{thm:certified}) are hopeless for dense networks.} On a 784--100--10 MNIST network (test error 2.4\%) we quantized the first-layer update in the DCT basis and certified with per-unit, per-band $\ell_2$ balls around the quantized network that contain $f$. At every step size the certified first-layer radius is 12--13 times the actual first-layer error $|Rx|$. Certifying 95\% of the training points needs steps of 2\% of the update's RMS, i.e.\ $6.1\times10^5$ bits for the first layer alone, ten times more than $m=6\times10^4$ (\repo{results/mnist/t3_worstcase_probe.txt}).

\paragraph{Dithered compression (\cref{thm:dither,lem:cert}).} \Cref{tab:t3} reports the bound for one-hidden-layer networks trained in the standard way (Adam, 20 epochs). It is non-vacuous (below the chance level of 0.9) for narrow networks: 0.44 for $H=8$ (test error 7.0\%) and 0.55 for $H=16$ (5.1\%) on MNIST, and 0.66 for $H=16$ on Fashion-MNIST. It degrades with width: 0.70 at $H=32$, 0.87 at $H=64$, 0.97 at $H=128$. Dithering does what it was designed to do. At comparable certified fractions it needs about 3 bits per first-layer coefficient instead of the 6--8 bits of worst-case sets (\repo{results/mnist/t3_worstcase_probe.txt}), because the certificate scales with the $\ell_2$ aggregate of the noise rather than its $\ell_1$ worst case. But at the measured 2.1--3.6 bits per first-layer coefficient (3.2--3.6 for $H\le32$), $784H$ coefficients exceed $m=6\times10^4$ bits once $H\gtrsim25$, and the total grows linearly in $H$.

Three ablations on the 32-unit network: coding the first layer in the pixel basis instead of the DCT basis costs 36\% more bits (bound 0.81 instead of 0.70); training with Gaussian weight noise of one or two times the initial RMS changes little (0.71, 0.69); a deterministic DCT initialization gives a more compressible network (67k instead of 84k bits; bound 0.69).

\paragraph{Networks with sparse updates.} The cost is driven by the number of first-layer coefficients that have to be described. Since exactly-zero coefficients are sent exactly (no dither, no noise, \cref{sec:methods}), a network whose update is sparse in the DCT basis is much cheaper. We trained 100- and 256-unit networks normally, kept the largest 0.1--3\% of the update's DCT coefficients and fine-tuned with that support fixed. Pruning followed by fine-tuning is a standard way to obtain small deployable networks \cite{zhou2019deconstructing}, and the bound is still on the network that was trained and deployed. \Cref{tab:t3sparse} shows the result. A standard-width MNIST network (100 hidden units) with 1\% of its update coefficients nonzero has test error 4.6\% and a training-data-only bound of 0.31. Keeping 0.3\% gives 7.2\% and 0.30, and 256-unit networks reach 0.40. The description is then 11--32 thousand bits. Each kept first-layer coefficient costs 5--8 bits for its value and 6--9 bits for its position, and the output layer takes 20--75\% of the total. For the output layer the search always picked the largest step multiplier on its grid, so a wider grid would probably help a little.

\begin{table}[t]
\centering\small
\begin{tabular}{l r l l r r r r r}
\toprule
data & $H$ & training & code & test err.\ $f$ & bits & certified & bound on $f$ & 32-bit Occam \\
\midrule
mnist & 32 & standard & DCT & 0.0323 & 84106 & 0.950 & 0.703 & 1.0 \\
mnist & 64 & standard & DCT & 0.0250 & 158075 & 0.966 & 0.871 & 1.0 \\
\bottomrule
\end{tabular}
\caption{Bounds on the original one-hidden-layer network $f$ from the 60{,}000 training examples alone (\cref{thm:dither,lem:cert}): no validation set and no unlabeled data. `Certified': fraction of training points with $\pi_g=1$. `Noise $\sigma$': $f$ trained with Gaussian weight noise of $\sigma$ times each tensor's initial RMS. Chance level is 0.9.}\label{tab:t3}
\end{table}


\begin{table}[t]
\centering\small
\begin{tabular}{l r r r r r r r}
\toprule
data & $H$ & kept coefficients & test err.\ $f$ & bits & certified & bound on $f$ & 32-bit Occam \\
\midrule
mnist & 100 & 3.0\% & 0.0333 & 32455 & 0.963 & 0.407 & 1.00 \\
mnist & 100 & 1.0\% & 0.0462 & 17664 & 0.945 & 0.314 & 1.00 \\
mnist & 100 & 0.3\% & 0.0722 & 11423 & 0.909 & 0.304 & 1.00 \\
mnist & 256 & 0.3\% & 0.0441 & 23244 & 0.924 & 0.399 & 1.00 \\
mnist & 256 & 0.1\% & 0.0647 & 21673 & 0.916 & 0.398 & 1.00 \\
fashion & 100 & 1.0\% & 0.1693 & 20176 & 0.828 & 0.501 & 1.00 \\
fashion & 100 & 0.3\% & 0.1926 & 13170 & 0.798 & 0.468 & 1.00 \\
\bottomrule
\end{tabular}
\caption{As \cref{tab:t3}, for networks whose first-layer update is exactly sparse in the DCT basis (trained, pruned to the stated fraction of coefficients, fine-tuned with the support fixed). The support is part of the description; only nonzero coefficients are dithered.}\label{tab:t3sparse}
\end{table}


\paragraph{Boolean networks: a bound for every input distribution.} \Cref{tab:boolean_t3} applies the same theorem to one-hidden-layer networks on Boolean tasks. Since it uses only the training sample, the bound holds whatever distribution generated the inputs. The bounds are 0.06--0.33 for networks with 16 or 32 hidden units at $m=16000$ (0.23 at $m=4000$). With the uniform distribution known, \cref{thm:reduction} on the same networks is 30--150 times tighter (0.0015--0.007), which measures what knowledge of the input distribution buys.

\begin{table}[t]
\centering\small
\begin{tabular}{l r r r r r r r r}
\toprule
task & $n$ & $m$ & $H$ & test err.\ $f$ & bits & training-data-only bound & bits & known-$D_X$ bound \\
\midrule
\texttt{maj5} & 20 & 4000 & 16 & 0.0000 & 1529 & 0.234 & 38 & 0.007 \\
\texttt{maj5} & 20 & 16000 & 16 & 0.0000 & 1523 & 0.064 & 38 & 0.002 \\
\texttt{parity3} & 20 & 16000 & 16 & 0.0000 & 2009 & 0.084 & 31 & 0.002 \\
\texttt{and3} & 50 & 16000 & 16 & 0.0000 & 3791 & 0.152 & 32 & 0.002 \\
\texttt{fullmaj} & 50 & 16000 & 16 & 0.0000 & 3629 & 0.151 & 48 & 0.002 \\
\texttt{junta6} & 50 & 16000 & 32 & 0.0000 & 9132 & 0.327 & 102 & 0.005 \\
\texttt{maj5} & 50 & 16000 & 16 & 0.0000 & 2589 & 0.110 & 45 & 0.002 \\
\texttt{maj5}, noise 1 & 50 & 16000 & 16 & 0.0000 & 2331 & 0.099 & 45 & 0.002 \\
\texttt{maj5} & 50 & 16000 & 32 & 0.0000 & 5468 & 0.213 & 45 & 0.002 \\
\texttt{parity3} & 50 & 16000 & 32 & 0.0000 & 7113 & 0.266 & 35 & 0.002 \\
\bottomrule
\end{tabular}
\caption{Boolean tasks, one-hidden-layer networks: the training-data-only bound (\cref{thm:dither}), which holds whatever the input distribution, against the bound that uses the known uniform input distribution (\cref{thm:reduction}) on the same network.}\label{tab:boolean_t3}
\end{table}

